\documentclass[10pt,a4paper]{amsart}
\usepackage[margin=19mm]{geometry}
\usepackage{amsmath,amssymb,mathtools}
\usepackage[scr=rsfs,cal=euler]{mathalfa}
\usepackage{xfrac}
\usepackage[unicode,hidelinks,bookmarks=false]{hyperref}
\newcommand{\ie}{i.e.\ }
\newcommand{\cf}{cf.\ }
\newcommand{\resp}{respectively\ }
\newcommand{\Subsection}{\subsection}
\providecommand{\nobreakdash}{\nobreak\hbox{-}}
\setlength{\emergencystretch}{2em}
\multlinegap0pt
\usepackage{bm}
\usepackage{bbm}
\usepackage{dsfont}
\usepackage{xspace}
\usepackage{cancel}
\usepackage{mathtools}
\usepackage{amsthm}
\makeatletter
\@ifundefined{theorem}{\newtheorem{theorem}{Theorem}}{}
\@ifundefined{lemma}{\newtheorem{lemma}[theorem]{Lemma}}{}
\makeatother
\newcommand{\Lam}{\Lambda}
\newcommand{\Star}{*}
\title[The Mantovanelli-Hofstadter Sequence]{The Mantovanelli-Hofstadter Sequence}
\author{Benoit Cloitre}
\date{Source: arXiv:2604.06237v3 (CC BY 4.0)}
\begin{document}
\maketitle

\noindent\emph{Source and licence.} The Mantovanelli-Hofstadter Sequence. Version arXiv:2604.06237v3 (CC BY 4.0), distributed under the Creative Commons Attribution 4.0 International licence, as recorded in the arXiv licence metadata for this exact version and shown on the abstract page https://arxiv.org/abs/2604.06237v3. Full rights notice: licenses/arxiv-2604.06237-CC-BY-4.0.txt.\par\smallskip

\noindent\textit{Editorial scope.} Counted unit: the Theorem whose source label is \texttt{thm:root-classification} in the article (source0.tex lines 1749--1762, source label \texttt{thm:root-classification}), delivered together with the article's own complete original proof of it (source0.tex lines 1764--1926, one \texttt{proof} environment of 163 lines). No mathematics is written, repaired or re-derived: statement and proof are the article's own text line for line. Required context retained uncounted: eq:phi-binomial-length (lines 1639--1642); eq:run-vector (lines 1651--1655); lem:root-suspension (lines 1668--1680). Every definition, display and label of those lines is retained unchanged. Omitted from this package and not redistributed: 1--1638, 1643--1650, 1656--1667, 1681--1748, 1763--1763, 1927--4485. Nothing inside a retained excerpt is omitted or altered: every retained body is delivered exactly as the article's lines are written, so no omission receipt of any kind accompanies this package. Citations: the retained excerpts carry no citation command at all, so no bibliography entry of the article is transcribed and no reference is redistributed. Reference adaptation: none was needed and none was made; every reference inside the delivered unit resolves inside it, so no unresolved reference or citation is produced. Printed slips kept literal: no printed word, symbol, hypothesis or number of the counted statement or of its proof was altered. Layout and class: the article's own class is \texttt{article} with packages including geometry, fontenc, inputenc, lmodern, microtype, amsmath, amssymb, amsthm, mathtools, mathrsfs, graphicx, tikz, enumitem, needspace; the wrapper is the lane's pinned \texttt{amsart} editorial wrapper at 10pt on A4, which restates the theorem-like environments the retained text needs and adds the article's own macro definitions that the retained text needs (\texttt{\textbackslash Lam}, \texttt{\textbackslash Star}), with extra wrapper packages (bm, bbm, dsfont, xspace, cancel, mathtools). The delivered numbering is the unit's own. Assets: no retained span contains a figure environment, a graphics-inclusion command, a PDF-page inclusion, a table or a TikZ picture; no image file is copied into this package, the package declares no asset dependency and no image file is redistributed with it. Licence: version arXiv:2604.06237v3 of this article is distributed under the Creative Commons Attribution 4.0 International licence, as recorded in the arXiv licence metadata for this exact version and shown on the abstract page https://arxiv.org/abs/2604.06237v3. A full rights notice is delivered at licenses/arxiv-2604.06237-CC-BY-4.0.txt. Characters: every retained excerpt is pure ASCII, so the delivered engine, XeLaTeX with its default Latin Modern text font, reports no missing character.
\begin{equation}
\label{eq:phi-binomial-length}
 |\varphi^k(L)|=\binom Lk
\end{equation}

\begin{equation}
\label{eq:run-vector}
 (j_1-1,c_{j_1},j_2-j_1,c_{j_2},\ldots,
 c_{j_s},n-j_s+1,\varrho).
\end{equation}

\begin{lemma}
\label{lem:root-suspension}
Let \(U\) be the canonical word of a nonempty root word \(R\).  Then the
canonical word of
\[
 \Gamma(R)=(1,R+2)
\]
is
\[
 \Lam(U)=T^2(U)[1:|T^2(U)|-1).
\]
The map \(\Lam\) is injective and commutes with \(\Star\).
\end{lemma}

\begin{theorem}
\label{thm:root-classification}
The roots of the nonempty anti-palindromic Dyck words satisfying
\(\mathsf S\) are exactly
\begin{equation*}
 R=0^b\qquad(b\geq1)
\end{equation*}
and
\begin{equation}
\label{eq:classified-root-positive}
 R=\mathcal J_M(M-1)^aM^b
 \qquad(M\geq1,\ a\geq0,\ b\geq1).
\end{equation}
\end{theorem}

\begin{proof}
The cases \(M=0,1,2\) follow directly from the run vector.  For
\(M=0\), one has \(R=0^b\).  Suppose \(M=1\), let \(b\) be the number of
letters one in \(R\), and let \(z\) be the number of its initial zeros.
Since \(c=0^bR\), the initial zero run has length \(b+z\), whereas the
terminal one run has length \(|R|\).  Their equality forces every zero of
\(R\) to occur before every one.  Thus \(R=0^a1^b\).

Suppose \(M=2\).  Let \(a\) count the letters one before the first letter two
of \(R\), and let \(b\) count all letters two.  The identities
\[
 \varphi(1)=(0),
 \qquad \varphi(2)=(1,0),
 \qquad \varphi^2(2)=(0)
\]
show that the initial zero run has length \(a+b\).  Equality with the
terminal run \(|R|\) excludes both a zero anywhere in \(R\) and a one after
the first two.  Hence \(R=1^a2^b\).

Assume \(M\geq3\).  Equation \eqref{eq:phi-binomial-length} shows that
\(\varphi^{M-1}(M)\) has length \(M\) and contains only zeros and ones.  For
every word \(w\) of nonnegative integers,
\[
 |\varphi(w)|=\sum_i w_i.
\]
Since \(|\varphi^M(M)|=1\), the sum of the letters of
\(\varphi^{M-1}(M)\) is one.  Thus this word contains one letter one and
\(M-1\) zeros.

Let \(h_M\) be the number of zeros before that one.  In \(\varphi(M)\), after
another \(M-2\) iterations, only the letter \(M-1\) contributes the unique
one.  The letter \(M-2\) contributes one zero, and every smaller letter has
disappeared.  The explicit definition of \(\sigma_M\) places \(M-2\) before
\(M-1\) in \(\varphi(M)\) exactly when \(M\) is odd.  Hence
\[
 h_1=0,
 \qquad
 h_M=h_{M-1}+\mathbf 1_{\{M\text{ is odd}\}}.
\]
Therefore \(h_M=\lfloor(M-1)/2\rfloor\), and
\begin{equation*}
 \varphi^M(M)=(0),
 \qquad
 \varphi^{M-1}(M)=0^t1\,0^{M-t-1},
 \qquad t=\left\lfloor\frac{M-1}{2}\right\rfloor.
\end{equation*}
Let \(b\) be the number of root letters \(M\), and let \(a\) be the number
of root letters \(M-1\) before the first \(M\).  The first zero run in the
degree word has length \(a+b+t\).  Anti-palindromicity gives
\begin{equation*}
 \varrho=|R|=a+b+t.
\end{equation*}

The positive nonroot degrees cannot end inside this initial comparison.  Let
\(\nu_x\) be the number of positive-degree nonroot descendants below a root of
type \(x\).  The root has one child of each type from zero through \(x-1\),
so
\[
 \nu_1=0,
 \qquad
 \nu_x=(x-1)+\sum_{y=1}^{x-1}\nu_y.
\]
Induction gives \(\nu_x=2^{x-1}-1\).  Let \(N_+(w)\) denote the number of
positive letters of an integer word \(w\).  Consequently,
\[
 N_+(\varphi(c))\geq b\nu_M+a\nu_{M-1}.
\]
Subtracting \(\varrho=a+b+t\) yields
\begin{equation*}
 N_+(\varphi(c))-\varrho
 \geq b(2^{M-1}-2)+a(2^{M-2}-2)-t\geq1.
\end{equation*}
For \(M=3\), the last expression is at least \(2b-1\).  For \(M\geq4\), the
term containing \(b\) suffices.  This is a lower bound, not an equality.
In particular, the first \(\varrho+1\) positive letters of \(c\) lie in the
prefix \(\varphi(c)=c_1\cdots c_{n-\varrho}\), strictly before \(R\).

Write
\[
 v_i=c_{j_i}>0,
 \qquad
 z_i=j_{i+1}-j_i\quad(1\leq i<s),
 \qquad
 z_s=n-j_s+1.
\]
The word \(c\) begins with \(0^{\varrho}\).  In the factorization
\(c=\varphi(c)R\), only a factor \(\varphi(1)=(0)\) can produce this initial
zero run, because \(\varphi(L)\) begins with one for \(L\geq2\).  There are
therefore exactly \(\varrho\) source letters one before the first source
letter at least two.  The preceding margin places their positive outputs
strictly before the root suffix.  Consequently,
\[
 v_1=\cdots=v_{\varrho}=1,
 \qquad v_{\varrho+1}\geq2.
\]
For \(i<\varrho\), the output \(v_i=1\) must be the unique positive output of
\(\varphi(2)=(1,0)\).  Otherwise it would be the first output of a factor
\(\varphi(L)\) with \(L\geq3\), whose next positive output is at least two,
contrary to \(v_{i+1}=1\).  Hence \(z_i\geq2\).  For \(i=\varrho\), the
factor cannot be \(\varphi(2)\), since the next positive output would again
begin with one, contrary to \(v_{\varrho+1}\geq2\).  It begins a factor
\(\varphi(L)\) with \(L\geq3\), whose first two positive outputs are adjacent.
Thus \(z_{\varrho}=1\).

Palindromicity of \eqref{eq:run-vector} gives, for
\(1\leq i\leq\varrho\), the two families of equalities
\[
 z_{s+1-i}=v_i=1,
 \qquad
 v_{s+1-i}=z_i.
\]
The first family says
\(j_s=n,j_{s-1}=n-1,\ldots,j_{s+1-\varrho}=n+1-\varrho\).  Thus the last
\(\varrho\) degree positions are positive and form the root word.  The
second family gives
\[
 R_{\varrho+1-i}=z_i
 \qquad(1\leq i\leq\varrho).
\]
Consequently,
\begin{equation}
\label{eq:root-descent}
 R=(1,R'+2)
\end{equation}
for a root word \(R'\) of maximum \(M-2\).  It is nonempty, because
\(M\geq3\) gives \(t\geq1\), and hence
\(\varrho=a+b+t\geq b+t\geq2\).

Let \(U\) be the canonical word of \(R'\).  Lemma
\ref{lem:root-suspension} identifies the canonical word of \(R\) with
\(\Lam(U)\).  Since this word is fixed by \(\Star\), injectivity and
commutation give \(U=U^{\Star}\).  Thus \(R'\) is again an anti-\(\mathsf S\)
root.  Iterating \eqref{eq:root-descent} reaches a base of maximum zero, one,
or two.  The identity
\[
 \mathcal J_M=(1,\mathcal J_{M-2}+2)
\]
then gives \eqref{eq:classified-root-positive}.

Conversely, the three base families have canonical words
\[
 0^b1^b,
 \qquad
 0^{a+b}(10)^b1^{a+b},
\]
and
\[
 0^{a+b}(100)^b(10)^a(110)^b1^{a+b}.
\]
Their outer runs are interchanged by \(\Star\), while
\[
 ((10)^b)^{\Star}=(10)^b
\]
and
\[
 \bigl((100)^b(10)^a(110)^b\bigr)^{\Star}
 =(100)^b(10)^a(110)^b.
\]
Thus the extreme runs \(0^{a+b}\) and \(1^{a+b}\) correspond exactly, the
three words satisfy \(\mathsf S\), and they are fixed by \(\Star\).  Lemma
\ref{lem:root-suspension} preserves both properties.  This proves the
converse and completes the classification.
\end{proof}

\end{document}
